%%%%%%%%%%%%%%%%%%%%%%%%%%%%%%%%%%%%%%%%%%%%%%%%%%
%%%%%%%%%%%%%%%%%%%%%%%%%%%%%%%%%%%%%%%%%%%%%%%%%%

% A primeira versão (1.0) deste modelo foi desenvolvida por Mauro Moura Severino, com a ajuda de Rubens Vasconcellos Terra Neto, para o Mestrado Profissional em Poder Legislativo (MPPL) da Câmara dos Deputados a partir do Modelo Canônico de TCC com ABNTEX2, elaborado pelo "abnTeX2 group". Cada uma das demais versões (2.0, 2.1 e 3.0) foi implementada por Mauro Moura Severino a partir da versão anterior.

%%%%%%%%%%%%%%%%%%%%%%%%%%%%%%%%%%%%%%%%%%%%%%%%%%

% Template de DISSERTAÇÃO
% Versão: v-3.0
% Data: 31/12/2024

% Esta versão viabiliza a produção de documentos com grau muito elevado de conformidade com as seguintes normas:
    % ABNT NBR 6023:2018 – Informação e documentação – Referências – Elaboração
    % ABNT NBR 6024:2012 – Informação e documentação – Numeração progressiva das seções de um documento – Apresentação
    % ABNT NBR 6027:2012 – Informação e documentação – Sumário – Apresentação
    % ABNT NBR 6028:2021 – Informação e documentação – Resumo, resenha e recensão – Apresentação
    % ABNT NBR 6034:2004 – Informação e documentação – Índice – Apresentação
    % ABNT NBR 10520:2023 – Informação e documentação – Citações em documentos – Apresentação
    % ABNT NBR 14724:2024 – Informação e documentação – Trabalhos acadêmicos – Apresentação
    % IBGE. Normas de apresentação tabular. 3. ed. Rio de Janeiro: IBGE, 1993.

%%%%%%%%%%%%%%%%%%%%%%%%%%%%%%%%%%%%%%%%%%%%%%%%%%
%%%%%%%%%%%%%%%%%%%%%%%%%%%%%%%%%%%%%%%%%%%%%%%%%%

% A EDIÇÃO DESTE ARQUIVO É EXCLUSIVA DA EQUIPE DA COORDENAÇÃO DE PÓS-GRADUAÇÃO DO CEFOR. O AUTOR NÃO DEVE EDITÁ-LO.

%%%%%%%%%%%%%%%%%%%%%%%%%%%%%%%%%%%%%%%%%%%%%%%%%%
%%%%%%%%%%%%%%%%%%%%%%%%%%%%%%%%%%%%%%%%%%%%%%%%%%

Este arquivo tem como objetivo orientar, de modo sucinto, a produção de um TCC para o MPPL, na modalidade DISSERTAÇÃO, a partir deste projeto do Overleaf.

1) O projeto está estruturado conforme pode ser visto no menu à esquerda. ESSA ESTRUTURAÇÃO NÃO DEVE SER ALTERADA.

2) O arquivo TCC-MPPL-Cefor-modelo-DIS.tex é a base do modelo, sendo o arquivo .tex a ser compilado para que haja o resultado no formato .pdf.

3) O autor deve editar APENAS os seguintes arquivos, conforme as instruções que constam em cada um:

3.1) na pasta 0-preâmbulo: I-opções.tex, II-autoria.tex e III-abreviaturas.tex;
3.2) na pasta 1-pré-textuais: I-dedicatória.tex, II-agradecimentos.tex, III-epígrafe.tex e IV-resumos.tex;
3.3) na pasta 2-textuais: I-introdução.tex, II-desenvolvimento.tex e III-conclusão.tex;
3.4) na pasta 3-pós-textuais: I-glossário.tex, II-apêndices.tex e III-anexos.tex; e
3.5) na pasta B-referências: TCC-MPPL-referências.bib.

4) Na pasta A-figuras, devem ser incluídos os arquivos de todas as figuras a serem utilizadas.

5) Os arquivos abntex2-alf.bst, abntex2.cls, abntex2cite.sty e TCC-MPPL-Cefor-modelo-DIS.tex NÃO DEVEM SER EDITADOS PELOS AUTORES.

6) Para melhor acompanhamento da construção do TCC, sugere-se que o autor faça a compilação do arquivo TCC-MPPL-Cefor-modelo-DIS.tex logo após cada realização de pequenos trechos de edição. Essa estratégia facilita a localização de eventuais erros e melhora o entendimento do funcionamento do LaTeX neste modelo do Overleaf.

7) A equipe da Copos estará à disposição do autor para ajudá-lo no que for necessário e conta com as contribuições dele para melhorar este modelo.